% Lecture 2 notes: \input by notes.tex, \subimport'ed by the semester book.
% Exam questions (зачёт): 1 (сложение, умножение), 5 (Карацуба), 6 (Тоом–Кук).
\providecommand{\Len}{\mathcal{L}}
\begin{lecture}{2}{Арифметика больших чисел}

\section{Зачем нужна «длинная» арифметика}

Криптографические протоколы, которые мы будем изучать (Диффи"--~Хеллман, RSA, протоколы на
эллиптических кривых), работают с числами длиной в сотни и тысячи двоичных разрядов: модуль RSA
сегодня "--- не меньше $2048$ бит, т.\,е. около $617$ десятичных цифр. Процессор умеет
непосредственно складывать и умножать только \emph{машинные слова} (32 или 64 бита), поэтому
операции над большими числами приходится программировать: число хранится как массив «цифр»
по основанию $B = 2^{32}$ или $B = 2^{64}$, а сложение, умножение и деление выражаются через
операции над цифрами.

Стойкость протоколов опирается на \emph{сложные} задачи (дискретное логарифмирование,
факторизация), а «законные» операции "--- шифрование, подпись, согласование ключа "--- сводятся
к умножениям и возведениям в степень больших чисел. Поэтому скорость этих операций критична:
в этой лекции мы увидим, что умножение $N$-значных чисел можно выполнить существенно быстрее,
чем за $N^2$ операций над цифрами.

\section{Представление больших чисел}

Всюду $B \ge 2$ "--- \emph{основание} системы счисления, элементы множества
$\{0, 1, \dots, B - 1\}$ называются \emph{цифрами}.

\begin{definition}
\emph{Представлением} натурального числа $x$ по основанию $B$ называется кратчайшая
последовательность цифр $(x_0, x_1, \dots, x_{D-1})$, для которой
\[
  x = \sum_{i=0}^{D-1} x_i B^i, \qquad 0 \le x_i < B.
\]
Пишут $x = (x_{D-1} \dots x_1 x_0)_B$. Число $D$ цифр называется \emph{$B$-длиной} числа $x$
и обозначается $\Len(x) = \Len_B(x)$; по определению $\Len(0) = 1$.
\end{definition}

Требование «кратчайшая» запрещает ведущие нули: $x_{D-1} \ne 0$ при $x > 0$.

\begin{theorem}\label{l2:repr}
Каждое натуральное $x$ имеет единственное представление по основанию $B$, и
$\Len(x) = \floor{\log_B x} + 1$ при $x \ge 1$.
\end{theorem}

\begin{proof}
\emph{Существование} "--- индукция по $x$: при $x < B$ представление $(x)$ состоит из одной
цифры. При $x \ge B$ разделим с остатком: $x = qB + r$, $0 \le r < B$, $1 \le q < x$; по
предположению индукции $q = \sum_{i=0}^{D-2} q_i B^i$, и тогда
$x = \sum_{i=0}^{D-2} q_i B^{i+1} + r$, т.\,е. $(q_{D-2}, \dots, q_0, r)$ "--- представление $x$.

\emph{Единственность}. Если $\sum x_i B^i = \sum y_i B^i$, то, приводя по модулю $B$, получаем
$x_0 \equiv y_0 \pmod B$, а так как $0 \le x_0, y_0 < B$, то $x_0 = y_0$. Вычитая и деля на $B$,
приходим к тому же равенству для более коротких последовательностей.

\emph{Длина}. Если $x_{D-1} \ne 0$, то
$B^{D-1} \le x \le \sum_{i=0}^{D-1} (B-1) B^i = B^D - 1 < B^D$, т.\,е.
$D - 1 \le \log_B x < D$, откуда $D = \floor{\log_B x} + 1$.
\end{proof}

Доказательство существования одновременно даёт алгоритм перевода в систему с основанием $B$:
цифры "--- последовательные остатки от деления на $B$, начиная с младшей.

\subsection{Сбалансированное представление}

Пусть $B$ чётно. Вместо цифр $0, \dots, B - 1$ можно использовать цифры из симметричного
диапазона
\[
  -\tfrac{B}{2} \le \hat x_i < \tfrac{B}{2}.
\]

\begin{definition}
\emph{Сбалансированным представлением} целого числа $x$ по чётному основанию $B$ называется
кратчайшая последовательность $(\hat x_0, \dots, \hat x_{D-1})$ таких цифр, что
$x = \sum_{i=0}^{D-1} \hat x_i B^i$.
\end{definition}

\begin{proposition}\label{l2:balanced}
Пусть $B \ge 4$ чётно. Каждое целое число (в том числе отрицательное) имеет единственное
сбалансированное представление.
\end{proposition}

\begin{proof}
Цифры $-B/2, \dots, B/2 - 1$ образуют полную систему вычетов по модулю $B$, поэтому младшая
цифра $\hat x_0 \equiv x \pmod B$ определена однозначно, и далее рассуждаем, как
в теореме~\ref{l2:repr}, для числа $x' = (x - \hat x_0)/B$. Процесс конечен: при $x \ne 0$
\[
  \abs{x'} \le \frac{\abs{x} + B/2}{B} < \abs{x},
\]
так как последнее неравенство равносильно $B/2 < (B - 1)\abs{x}$, что верно при $\abs{x} \ge 1$
и $B \ge 4$. Значит, $\abs{x}$ строго убывает, пока не станет равным нулю.
\end{proof}

\begin{example}
Перевод из обычного представления в сбалансированное выполняется одним проходом от младших
разрядов: если $x_i + c \ge B/2$ ($c$ "--- перенос), записываем цифру $x_i + c - B$ и переносим
единицу в следующий разряд. Для $B = 10$:
\[
  1979 = (1\,9\,7\,9)_{10}: \quad 9 \to -1\ (\text{перенос } 1),\ 7 + 1 \to -2\ (1),\
  9 + 1 \to 0\ (1),\ 1 + 1 \to 2,
\]
т.\,е. $1979 = 2 \cdot 10^3 + 0 \cdot 10^2 - 2 \cdot 10 - 1$.
\end{example}

Сбалансированное представление удобно тем, что знак числа хранится «в цифрах», а суммы
произведений цифр (свёртки, см. лекцию~4) оказываются по модулю вдвое меньше, чем при обычных
цифрах.

\checkyourself{Как выглядит сбалансированное представление числа $-1979$ по основанию $10$?
Почему утверждение~\ref{l2:balanced} неверно при $B = 2$? Почему для нечётного $B$ естественно
брать цифры $-\frac{B-1}{2}, \dots, \frac{B-1}{2}$?}

\section{Оценки сложности}

Сложность алгоритмов длинной арифметики измеряется числом \emph{элементарных операций} "---
операций над цифрами (сложение и умножение двух цифр с получением двухцифрового результата,
деление двухцифрового числа на цифру) "--- как функция длины входа $N$.

Пусть $g(x) > 0$ при всех достаточно больших $x$.

\begin{definition}
\begin{enumerate}
  \item $f(x) = \bigO(g(x))$ при $x \to \infty$, если существуют $C > 0$ и $x_C$ такие, что
    $\abs{f(x)} \le C g(x)$ при всех $x > x_C$. Аналогично при $x \to a$: неравенство
    выполняется при $0 < \abs{x - a} < \delta_C$, т.\,е. $\limsup_{x \to a} \abs{f(x)}/g(x) < \infty$.
  \item $f(x) = o(g(x))$, если для любого $\varepsilon > 0$ найдётся $N_\varepsilon$:
    $\abs{f(x)} \le \varepsilon g(x)$ при $x > N_\varepsilon$, т.\,е. $\lim f(x)/g(x) = 0$.
  \item $f = \Omega(g)$, если $g = \bigO(f)$; $f = \Theta(g)$, если $f = \bigO(g)$ и $g = \bigO(f)$
    (обозначения Д.~Кнута).
  \item $f \sim g$ (\emph{асимптотическая эквивалентность}), если $\lim f(x)/g(x) = 1$.
\end{enumerate}
\end{definition}

Например, $3N^2 + 10N = \Theta(N^2)$, $N \log N = o(N^{1{,}01})$, $\Len_B(x) \sim \log_B x$.

\section{Сложение и вычитание}

Пусть $u = (u_{n-1} \dots u_0)_B$ и $v = (v_{n-1} \dots v_0)_B$ (более короткое число дополнено
ведущими нулями).

\begin{algorithm}[сложение, алгоритм A Кнута]
\begin{algorithmic}[1]
\Require $u, v$ "--- $n$-значные числа.
\Ensure $w = u + v = (w_n w_{n-1} \dots w_0)_B$.
\State $k \gets 0$ \Comment{перенос}
\For{$j = 0, 1, \dots, n - 1$}
  \State $t \gets u_j + v_j + k$
  \State $w_j \gets t \bmod B$, \quad $k \gets \floor{t / B}$
\EndFor
\State $w_n \gets k$
\end{algorithmic}
\end{algorithm}

\begin{lemma}
Перенос $k$ всегда равен $0$ или $1$.
\end{lemma}

\begin{proof}
Индукция по $j$: если $k \le 1$, то $t \le (B - 1) + (B - 1) + 1 = 2B - 1$, поэтому
$\floor{t/B} \le 1$.
\end{proof}

\begin{proposition}\label{l2:addlen}
$\max(\Len(a), \Len(b)) \le \Len(a + b) \le \max(\Len(a), \Len(b)) + 1$ для $a, b \in \NN$.
\end{proposition}

\begin{proof}
Пусть $D = \max(\Len(a), \Len(b))$. Тогда $B^{D-1} \le \max(a, b) \le a + b < 2B^D \le B^{D+1}$.
\end{proof}

Алгоритм выполняет $n$ шагов, каждый "--- $\bigO(1)$ операций над цифрами, т.\,е. сложность
сложения равна $\Theta(n)$. Это оптимально: любой алгоритм сложения должен хотя бы прочитать все
$2n$ цифр слагаемых. Вычитание $u - v$ при $u \ge v$ выполняется так же, с \emph{заёмом}
$k \in \{0, -1\}$: $t = u_j - v_j + k$, $w_j = t \bmod B$, $k = \floor{t/B}$.

\checkyourself{Сформулируйте алгоритм сложения двух чисел, записанных в сбалансированной форме.
Каким может быть перенос?}

\section{Школьное умножение}

Пусть $a = \sum_{i=0}^{M-1} a_i B^i$, $b = \sum_{j=0}^{N-1} b_j B^j$. Раскрывая скобки,
\[
  c = ab = \sum_{t=0}^{M+N-2} \Bigl(\sum_{i + j = t} a_i b_j\Bigr) B^t.
\]
Внутренние суммы $\sum_{i+j=t} a_i b_j$ "--- это \emph{линейная свёртка} последовательностей
цифр (подробно "--- в лекции~4); они могут превышать $B$, поэтому цифры произведения получаются
из них после переносов.

\begin{algorithm}[умножение, алгоритм M Кнута]
\begin{algorithmic}[1]
\Require $u = (u_{m-1} \dots u_0)_B$, $v = (v_{n-1} \dots v_0)_B$.
\Ensure $w = uv = (w_{m+n-1} \dots w_0)_B$.
\State $w_0 \gets \dots \gets w_{m-1} \gets 0$
\For{$j = 0, \dots, n - 1$}
  \State $k \gets 0$
  \For{$i = 0, \dots, m - 1$}
    \State $t \gets u_i v_j + w_{i+j} + k$
    \State $w_{i+j} \gets t \bmod B$, \quad $k \gets \floor{t / B}$
  \EndFor
  \State $w_{j+m} \gets k$
\EndFor
\end{algorithmic}
\end{algorithm}

\begin{lemma}
Во всех шагах алгоритма M $0 \le t < B^2$, т.\,е. $t$ помещается в два машинных слова,
а $k < B$.
\end{lemma}

\begin{proof}
Если $k \le B - 1$, то $t \le (B-1)^2 + (B-1) + (B-1) = B^2 - 1$.
\end{proof}

\begin{proposition}\label{l2:mullen}
$\Len(a) + \Len(b) - 1 \le \Len(ab) \le \Len(a) + \Len(b)$ при $a, b \ge 1$.
\end{proposition}

\begin{proof}
$B^{M-1} B^{N-1} \le ab < B^M B^N$, где $M = \Len(a)$, $N = \Len(b)$.
\end{proof}

Алгоритм M выполняет $mn$ умножений цифр и столько же сложений, т.\,е. умножение двух
$N$-значных чисел школьным методом требует $\Theta(N^2)$ операций.

\subsection{Возведение в квадрат}

При $b = a$ слагаемые $a_i a_{t-i}$ и $a_{t-i} a_i$ совпадают, поэтому
\[
  c_t = \sum_{i=0}^{t} a_i a_{t-i} = 2 \sum_{i=0}^{\floor{t/2}} a_i a_{t-i} - \mu_t, \qquad
  \mu_t = \begin{cases} 0, & t \text{ нечётно},\\ a_{t/2}^2, & t \text{ чётно}.\end{cases}
\]
(Слагаемое $i = t/2$ при чётном $t$ входит в сумму один раз, а удвоение учитывает его дважды.)
Различных произведений $a_i a_j$, $i \le j$, всего $N(N+1)/2$, поэтому возведение в квадрат
почти вдвое быстрее умножения.

\begin{proposition}
Если возведение в квадрат $N$-значного числа требует $S(N)$ операций, то умножение двух
$N$-значных чисел выполняется за $2S(N + 1) + \bigO(N)$ операций.
\end{proposition}

\begin{proof}
$ab = \frac14\bigl((a + b)^2 - (a - b)^2\bigr)$; числа $a \pm b$ имеют не более $N + 1$ цифр,
а сложения, вычитания и деление на $4$ (сдвиг) стоят $\bigO(N)$.
\end{proof}

Таким образом, умножение и возведение в квадрат имеют одинаковую \emph{асимптотическую}
сложность: умножение не может быть принципиально «сложнее».

\checkyourself{Почему в формуле для $c_t$ верхний предел суммы равен $\floor{t/2}$, а не
$\ceil{t/2}$? Проверьте формулу при $t = 1$ и $t = 2$.}

\section{Алгоритм Карацубы}

В 1960 г. А.~Н.~Колмогоров высказал гипотезу, что умножение $N$-значных чисел требует
$\Omega(N^2)$ операций. Через неделю его студент А.~А.~Карацуба опроверг её~\cite{karatsuba-ofman}.

Пусть $N = 2k$, $K = B^k$, и $a = a_1 K + a_0$, $b = b_1 K + b_0$, где $0 \le a_0, a_1, b_0, b_1 < K$
(т.\,е. мы записали $a$ и $b$ по основанию $K$). Тогда
\[
  ab = a_1 b_1 K^2 + (a_1 b_0 + a_0 b_1) K + a_0 b_0,
\]
и ключевое наблюдение состоит в том, что средний коэффициент выражается через \emph{одно}
дополнительное умножение:
\[
  a_1 b_0 + a_0 b_1 = (a_1 + a_0)(b_1 + b_0) - a_1 b_1 - a_0 b_0.
\]
Итак, умножение $N$-значных чисел сводится к трём умножениям чисел длины (примерно) $N/2$:
\[
  r_2 = a_1 b_1, \qquad r_0 = a_0 b_0, \qquad r_1 = (a_1 + a_0)(b_1 + b_0) - r_2 - r_0, \qquad
  ab = r_2 K^2 + r_1 K + r_0,
\]
плюс $\bigO(N)$ сложений, вычитаний и сдвигов (умножение на $K = B^k$ "--- сдвиг на $k$ цифр).

\begin{algorithm}[Карацуба]
\begin{algorithmic}[1]
\Function{Mul}{$a, b$} \Comment{$a, b < B^N$}
  \If{$N \le N_0$} \State \Return школьное произведение $ab$ \EndIf
  \State $k \gets \ceil{N/2}$, $K \gets B^k$; $a = a_1 K + a_0$, $b = b_1 K + b_0$
  \State $r_2 \gets \Call{Mul}{a_1, b_1}$, $r_0 \gets \Call{Mul}{a_0, b_0}$
  \State $r_1 \gets \Call{Mul}{a_1 + a_0, b_1 + b_0} - r_2 - r_0$
  \State \Return $r_2 K^2 + r_1 K + r_0$
\EndFunction
\end{algorithmic}
\end{algorithm}

\begin{theorem}\label{l2:karatsuba}
Алгоритм Карацубы умножает два $N$-значных числа за $\bigO(N^{\log_2 3})$ операций,
$\log_2 3 \approx 1{,}585$.
\end{theorem}

\begin{proof}
Пусть $T(N)$ "--- число операций. Суммы $a_1 + a_0$ и $b_1 + b_0$ могут иметь $k + 1$ цифру,
но умножение $(k+1)$-значных чисел сводится к умножению $k$-значных и $\bigO(k)$ дополнительным
операциям. Поэтому $T(N) \le 3T(\ceil{N/2}) + cN$ для некоторой константы $c$. Для $N = 2^m$,
раскрывая рекурсию,
\[
  T(2^m) \le 3^m T(1) + c \sum_{i=0}^{m-1} 3^i \, 2^{m-i}
  = 3^m T(1) + c\,2^m \frac{(3/2)^m - 1}{3/2 - 1} \le (T(1) + 2c)\, 3^m.
\]
Так как $3^m = (2^m)^{\log_2 3} = N^{\log_2 3}$, получаем $T(N) = \bigO(N^{\log_2 3})$; для
произвольного $N$ дополним числа нулями до ближайшей степени двойки, что увеличит $N$ не более
чем вдвое.
\end{proof}

\begin{remark}
Чтобы избежать «лишней» цифры у сумм, часто используют вариант с разностями:
$a_1 b_0 + a_0 b_1 = a_1 b_1 + a_0 b_0 - (a_1 - a_0)(b_1 - b_0)$; здесь $\abs{a_1 - a_0} < K$,
но приходится следить за знаком. На практике рекурсию останавливают не на однозначных числах,
а на длине $N_0$ порядка нескольких десятков машинных слов, где школьный метод быстрее
из-за меньших накладных расходов.
\end{remark}

\begin{example}\label{l2:karatsuba-hand}
$a = 12345$, $b = 678901$, $K = 10^4$: $a = 1 \cdot 10^4 + 2345$, $b = 67 \cdot 10^4 + 8901$.
\[
\begin{gathered}
  r_2 = 1 \cdot 67 = 67, \qquad r_0 = 2345 \cdot 8901 = 20872845,\\
  r_1 = 2346 \cdot 8968 - 67 - 20872845 = 21038928 - 20872912 = 166016,
\end{gathered}
\]
\[
  ab = 67 \cdot 10^8 + 166016 \cdot 10^4 + 20872845 = 8381032845.
\]
Умножения $2345 \cdot 8901$ и $2346 \cdot 8968$ можно, в свою очередь, выполнить тем же методом
с $K = 10^2$.
\end{example}

\taskexample{karatsuba}

\checkyourself{Сколько умножений однозначных чисел выполняет алгоритм Карацубы для двух
$2^m$-значных чисел (рекурсия до однозначных чисел)? Сравните с $4^m$ для школьного метода
при $m = 10$.}

\section{Основная теорема о рекуррентных соотношениях}

Оценки сложности алгоритмов «разделяй и властвуй» сводятся к рекуррентным соотношениям вида
$T(n) = a\,T(n/b) + f(n)$: задача размера $n$ разбивается на $a$ подзадач размера $n/b$, а
разбиение и сборка результата стоят $f(n)$. Число $c_{\mathrm{crit}} = \log_b a$ называется
\emph{критической экспонентой}.

\begin{theorem}[основная теорема, \cite{clrs}]\label{l2:master}
Пусть $a \ge 1$, $b > 1$, $T(n) = a\,T(n/b) + f(n)$.
\begin{enumerate}
  \item Если $f(n) = \bigO(n^c)$ для некоторого $c < c_{\mathrm{crit}}$, то
    $T(n) = \Theta(n^{c_{\mathrm{crit}}})$ (доминируют подзадачи).
  \item Если $f(n) = \Theta(n^{c_{\mathrm{crit}}} \log^k n)$, $k \ge 0$, то
    $T(n) = \Theta(n^{c_{\mathrm{crit}}} \log^{k+1} n)$ (вклад всех уровней одинаков).
  \item Если $f(n) = \Omega(n^c)$ для некоторого $c > c_{\mathrm{crit}}$ и выполнено
    \emph{условие регулярности} $a f(n/b) \le k f(n)$ для некоторого $k < 1$ и всех больших $n$,
    то $T(n) = \Theta(f(n))$ (доминирует разбиение).
\end{enumerate}
\end{theorem}

Докажем важный частный случай $f(n) = \Theta(n^c)$ для $n = b^m$ (общий случай "--- без
доказательства, см.~\cite[гл.~4]{clrs}).

\begin{proof}[Доказательство для $f(n) = \Theta(n^c)$]
На $i$-м уровне рекурсии $a^i$ подзадач размера $n/b^i$, их «собственная» стоимость
$a^i \Theta((n/b^i)^c) = \Theta(n^c q^i)$, где $q = a/b^c$. Суммируя по уровням $i = 0, \dots, m$,
\[
  T(n) = \Theta\Bigl(n^c \sum_{i=0}^{m} q^i\Bigr).
\]
Если $c > c_{\mathrm{crit}}$, то $q < 1$, сумма ограничена константой и $T = \Theta(n^c)$.
Если $c = c_{\mathrm{crit}}$, то $q = 1$ и $T = \Theta(n^c (m + 1)) = \Theta(n^c \log n)$.
Если $c < c_{\mathrm{crit}}$, то $q > 1$, сумма равна $\Theta(q^m)$, и
$T = \Theta(n^c a^m / b^{cm}) = \Theta(a^m) = \Theta(n^{\log_b a})$.
\end{proof}

\begin{example}
\begin{enumerate}
  \item $T(n) = 10\,T(n/3) + 42n^2$: $c = 2 = \log_3 9 < \log_3 10$, случай~1:
    $T = \Theta(n^{\log_3 10})$.
  \item $T(n) = 2\,T(n/4) + 23\sqrt n$: $c_{\mathrm{crit}} = \log_4 2 = \frac12 = c$, $k = 0$,
    случай~2: $T = \Theta(\sqrt n \log n)$.
  \item $T(n) = 4\,T(n/3) + n^2$: $c = 2 > \log_3 4$, и $4 (n/3)^2 = \frac49 n^2$, т.\,е.
    $k = \frac49 < 1$, случай~3: $T = \Theta(n^2)$.
  \item Карацуба: $T(N) = 3T(N/2) + \Theta(N)$, $c = 1 < \log_2 3$, случай~1.
\end{enumerate}
\end{example}

\checkyourself{К какому случаю основной теоремы относится школьное умножение, записанное
рекурсивно: $T(N) = 4T(N/2) + \Theta(N)$? Что даёт теорема?}

\section{Алгоритм Тоома"--~Кука}

Обобщим идею Карацубы: разобьём числа не на две, а на $D$ частей.

Пусть $N = Dk$, $K = B^k$, т.\,е. $a, b < K^D$. Запишем
\[
  a(x) = a_0 + a_1 x + \cdots + a_{D-1} x^{D-1}, \qquad
  b(x) = b_0 + b_1 x + \cdots + b_{D-1} x^{D-1},
\]
где $a_i, b_i$ "--- цифры чисел $a$ и $b$ по основанию $K$, так что $a(K) = a$, $b(K) = b$.
Например, $a = 126386$ при $K = 10$ даёт $a(x) = 6 + 8x + 3x^2 + 6x^3 + 2x^4 + x^5$, а при
$K = 100$ "--- $a(x) = 86 + 63x + 12x^2$.

\begin{lemma}\label{l2:eval}
Пусть $c(x) = a(x) b(x)$. Тогда $c(K) = ab$ и $\deg c \le 2D - 2$.
\end{lemma}

\begin{proof}
Подстановка $x = K$ "--- гомоморфизм колец $\ZZ[x] \to \ZZ$, поэтому $c(K) = a(K) b(K) = ab$.
Степени многочленов при умножении складываются.
\end{proof}

Таким образом, достаточно найти $2D - 1$ коэффициентов многочлена $c(x)$, после чего
$ab = c(K)$ вычисляется сдвигами и сложениями (с переносами, так как $c_k$ могут превышать $K$).
Коэффициенты находятся \emph{интерполяцией}.

\begin{theorem}\label{l2:interp}
Многочлен степени не выше $n$ однозначно определяется своими значениями в $n + 1$ различных
точках $x_0, \dots, x_n$.
\end{theorem}

\begin{proof}
Коэффициенты $c_0, \dots, c_n$ удовлетворяют линейной системе $\sum_k c_k x_i^k = c(x_i)$,
$i = 0, \dots, n$, с матрицей Вандермонда $V = (x_i^k)$. Её определитель равен
$\prod_{0 \le i < j \le n} (x_j - x_i) \ne 0$, поэтому система имеет единственное решение.
(Другое доказательство: разность двух таких многочленов имеет степень не выше $n$ и $n + 1$
корень, поэтому равна нулю.)
\end{proof}

\begin{algorithm}[Тоом"--~Кук]
\begin{enumerate}
  \item \textbf{[Подготовка]} Составить \emph{символические} многочлены $a(x)$, $b(x)$ степени
    $D - 1$ и выбрать $2D - 1$ точек, например $t \in \{1 - D, \dots, D - 1\}$.
  \item \textbf{[Вычисление]} Для каждой точки $t$ вычислить $a(t)$, $b(t)$ (сложения
    и умножения на малые константы) и \emph{одно длинное умножение} $c(t) = a(t) b(t)$
    чисел длины $\approx N/D$.
  \item \textbf{[Восстановление]} Решить систему $c(t) = \sum_{k=0}^{2D-2} c_k t^k$ относительно
    $c_k$. Решение выписывается заранее, \emph{символически}: каждое $c_k$ "--- фиксированная
    линейная комбинация значений $c(t)$ с рациональными коэффициентами.
  \item \textbf{[Формирование результата]} $ab = \sum_k c_k K^k$ (переносы).
\end{enumerate}
\end{algorithm}

Для $D = 3$ и точек $-2, -1, 0, 1, 2$ обозначим $t_0 = c(-2)$, $t_1 = c(-1)$, $t_2 = c(0)$,
$t_3 = c(1)$, $t_4 = c(2)$. Решая систему с матрицей Вандермонда, получаем
\[
\begin{aligned}
  c_0 &= t_2, &
  c_1 &= \tfrac{1}{12}t_0 - \tfrac23 t_1 + \tfrac23 t_3 - \tfrac1{12} t_4,\\
  c_2 &= -\tfrac1{24} t_0 + \tfrac23 t_1 - \tfrac54 t_2 + \tfrac23 t_3 - \tfrac1{24} t_4, &
  c_3 &= -\tfrac1{12} t_0 + \tfrac16 t_1 - \tfrac16 t_3 + \tfrac1{12} t_4,\\
  c_4 &= \tfrac1{24} t_0 - \tfrac16 t_1 + \tfrac14 t_2 - \tfrac16 t_3 + \tfrac1{24} t_4.
\end{aligned}
\]
Деления здесь \emph{точные}: значения $c_k$ "--- целые числа. На практике удобнее точки
$0, 1, -1, 2, \infty$, где «значение в бесконечности» "--- произведение старших коэффициентов
$c(\infty) = a_{D-1} b_{D-1} = c_{2D-2}$; тогда $c_0 = c(0)$, $c_4 = c(\infty)$,
$c_2 = \frac{c(1) + c(-1)}{2} - c_0 - c_4$, и остаётся система из двух уравнений на $c_1, c_3$
(см. пример~\ref{ex:toom-cook}).

\begin{example}
$a = 549238$, $b = 807541$, $K = 100$: $a(x) = 38 + 92x + 54x^2$, $b(x) = 41 + 75x + 80x^2$.
\[
\begin{aligned}
  t_0 &= (38 - 184 + 216)(41 - 150 + 320) = 70 \cdot 211 = 14770, &
  t_2 &= 38 \cdot 41 = 1558,\\
  t_1 &= (38 - 92 + 54)(41 - 75 + 80) = 0, & t_3 &= 184 \cdot 196 = 36064,\\
  t_4 &= 438 \cdot 511 = 223818.
\end{aligned}
\]
По формулам: $c_0 = 1558$, $c_1 = 6622$, $c_2 = 12154$, $c_3 = 11410$, $c_4 = 4320$, и
\[
  ab = 4320 \cdot 10^8 + 11410 \cdot 10^6 + 12154 \cdot 10^4 + 6622 \cdot 10^2 + 1558
     = 443532203758.
\]
\end{example}

\taskexample{toom-cook}

\begin{theorem}\label{l2:toom}
При фиксированном $D \ge 2$ алгоритм Тоома"--~Кука умножает $N$-значные числа за
$\Theta(N^{\log_D (2D - 1)})$ операций. Для любого $\varepsilon > 0$ существует $D$, при котором
сложность равна $\bigO(N^{1 + \varepsilon})$.
\end{theorem}

\begin{proof}
При фиксированном $D$ подготовка, восстановление (линейные комбинации с фиксированными
коэффициентами) и переносы стоят $\bigO(N)$, а длинных умножений чисел длины $N/D$ "---
$2D - 1$. Получаем $T(N) = (2D - 1)\,T(N/D) + \Theta(N)$, и по случаю~1 основной
теоремы~\ref{l2:master} ($1 < \log_D(2D - 1)$) $T(N) = \Theta(N^{\log_D(2D-1)})$. Так как
$\log_D (2D - 1) < \log_D (2D) = 1 + \log_D 2 \to 1$ при $D \to \infty$, для любого
$\varepsilon > 0$ найдётся $D$ с $\log_D (2D - 1) < 1 + \varepsilon$.
\end{proof}

Например, $D = 2$ даёт алгоритм Карацубы ($\log_2 3 \approx 1{,}585$), $D = 3$ "---
$\log_3 5 \approx 1{,}465$. Однако константа в $\bigO$ быстро растёт с $D$: число сложений
на этапах подготовки и восстановления растёт как $D^2$, а знаменатели коэффициентов
интерполяции "--- как факториал $(2D)!$ Поэтому на практике используют $D = 3$ (иногда $4$),
а для очень больших чисел "--- алгоритмы на основе быстрого преобразования Фурье.

\checkyourself{Почему в алгоритме Тоома"--~Кука нужно ровно $2D - 1$ точек? Что произойдёт,
если взять только $D$ точек?}

\checkyourself{Покажите, что при $D = 2$ и точках $0, 1, \infty$ алгоритм Тоома"--~Кука
совпадает с алгоритмом Карацубы.}

\section{Что дальше}

\begin{itemize}
  \item Алгоритм Шёнхаге"--~Штрассена~\cite{schonhage-strassen} (1971) умножает $N$-значные
    числа за $\bigO(N \log N \log\log N)$ операций с помощью быстрого преобразования Фурье
    (лекция~4). В 2019 г. Д.~Харви и Й.~ван~дер~Хувен построили алгоритм сложности
    $\bigO(N \log N)$.
  \item Деление с остатком сводится к умножению (лекция~3); возведение в степень "---
    к $\bigO(\log e)$ умножениям (лекция~6).
  \item В библиотеках (например, GMP) школьный метод, Карацуба, Тоом"--~Кук разных порядков
    и БПФ переключаются по порогам длины, подобранным экспериментально.
\end{itemize}

\checkyourselfsection

\begin{furtherreading}
  \readitem[разд.~4.3.1, 4.3.3]{knuth-taocp2}{школьные алгоритмы A, S, M; методы Карацубы и
    Тоома"--~Кука}
  \readitem[разд.~9.1, 9.5.1]{crandall-pomerance}{школьные методы, возведение в квадрат,
    Карацуба и Тоом"--~Кук}
  \readitem[гл.~4]{welschenbach}{реализация арифметики больших чисел на Си}
  \readitem[гл.~4]{clrs}{основная теорема о рекуррентных соотношениях и её доказательство}
  \readitem[разд.~14.2]{hac}{алгоритмы длинной арифметики}
  \readitem{karatsuba-ofman}{оригинальная статья}
\end{furtherreading}
\end{lecture}
